\chapter{От адронных реггевских траекторий к светоконусным струнам}

\lectureoverview{Теория струн началась не как теория квантовой гравитации.
Она возникла как попытка понять спектр и рассеяние адронов.
Прямые реггевские траектории указывали на протяжённый объект; дуальность в
пион--пионном рассеянии наводила на мысль, что два внешне различных процесса были
двумя разбиениями одной поверхности. Мы проследим эти указания до картины
мировой поверхности, переинтерпретируем адронную струну как трубку потока КХД,
а затем используем систему бесконечного импульса, чтобы вывести поперечную
механику открытой струны.
Главный результат состоит в том, что её эффективный внутренний гамильтониан измеряет
квадрат массы. Поэтому в этом описании энергия по закону Гука, пропорциональная \(L^2\),
превращается в энергию покоя, пропорциональную \(L\), — определяющее поведение
струны с постоянным натяжением.}

\section{Адронное происхождение теории струн}

Около 1968--1970 проблема, приведшая к струнам, касалась протонов,
нейтронов и особенно мезонов. Кварковая модель уже была серьёзной идеей,
но кварковый состав адронов ещё не был экспериментально установлен
в современном смысле. Глюоны в обычных обсуждениях фигурировали ещё меньше.
Намбу предлагал родственные идеи калибровочных полей, но устоявшегося языка
квантовой хромодинамики тогда ещё не существовало. Исходным материалом был
вместо этого неожиданно богатый спектр сильновзаимодействующих частиц.

Полезно временно пренебречь различием протон--нейтрон и называть
оба состояния нуклоном. Помимо нуклона были более тяжёлые состояния с большим
спином, за которыми следовали ещё более тяжёлые состояния с ещё большим
спином. Chew и Frautschi упорядочили такие семейства на графике, где угловой
момент откладывался по вертикальной оси, а квадрат массы — по горизонтальной.
В исторических единицах того времени масса протона составляла около одного BeV; теперь
та же единица называется GeV.

Будем измерять угловой момент в единицах \(\hbar\). Для барионного семейства
безразмерные значения спина полуцелые,

\begin{equation}
  J=\frac12,\frac32,\frac52,\ldots,
\end{equation}

тогда как мезонные семейства имеют целый спин. Примечательный эмпирический факт состоял
в том, что каждое семейство при построении графика по \(M^2\) располагалось приблизительно
на прямой. В современных обозначениях пишут

\begin{equation}
  J\simeq \alpha_0+\alpha' M^2,
  \label{eq:l01-regge-trajectory}
\end{equation}

где \(\alpha_0\) — перехват, а \(\alpha'\) — наклон реггевской траектории.
\editorialnote{В лекции подчёркиваются наблюдаемая линейная зависимость и её
почти универсальный наклон. Обозначение перехвата вводится только для того,
чтобы компактно записать эту зависимость; сами траектории эмпиричны и приблизительны.}

Изначальная прямая не выглядела убедительной, потому что проходила только
через две известные точки. Любые две точки задают прямую. Удивление пришло позже:
вновь открываемые состояния продолжали ложиться близко к той же прямой. Мезонные
траектории обнаруживали ту же организацию. Пионная траектория начинается вблизи \(J=0\) при
малой массе на шкале GeV; ро-траектория начинается при \(J=1\); дальнейшие возбуждения
продолжаются вверх. В одном семействе могло появляться пять, шесть или семь состояний. Ещё
более поразительно, что различные барионные и мезонные траектории были почти
параллельны. Тем самым увеличению углового момента на один квант соответствовало
приблизительно одно и то же увеличение квадрата массы:

\begin{equation}
  \Delta J=1
  \quad\Longrightarrow\quad
  \Delta M^2\simeq \frac{1}{\alpha'}.
  \label{eq:l01-regge-spacing}
\end{equation}

\begin{classroomqa}
\audiencequestion{Что означает, что разные реггевские траектории
параллельны?\lecturetimestamp{00:09:32}}
\lecturerresponse{Выберите одно семейство частиц и нанесите все его возбуждённые
состояния на график; они лежат вдоль одной прямой. Повторите это для другого мезонного или барионного
семейства. Прямые имеют приблизительно одинаковый наклон, хотя их
перехваты различаются. Поэтому одному и тому же приращению углового момента
соответствует приблизительно одно и то же приращение \(M^2\). Этот общий наклон называли
универсальным реггевским наклоном.}
\end{classroomqa}

Сравнение с обычными протяжёнными объектами показывает, почему эта закономерность имела значение.
Баскетбольный мяч приобретает энергию вращения, когда его раскручивают, но его траектория
скорее искривлена, чем линейна, и при достаточно большой угловой скорости мяч
разрывается. У атома тоже есть вращательные возбуждения, но в конце концов он
ионизуется; его траектория обрывается. Адроны демонстрировали гораздо более простую последовательность и,
в пределах имевшихся данных, не обнаруживали признаков того, что она исчерпывается.

Этот спектр говорил о том, что адрон ведёт себя не как точка. У точки нет
пространственной структуры, которую можно было бы раскручивать или деформировать. У электрона
не наблюдается башни внутренних состояний с большим спином, и, соответственно,
в пределах нынешних экспериментальных ограничений он точечен. Адрон же, напротив,
можно раскрутить. Его спектр указывает на составную структуру, а связь между
спином и массой намекает на то, что объект деформируется по мере роста его
углового момента.

\section{Дуальность в пион--пионном рассеянии}

Спектроскопия дала одну подсказку. Другую подсказку дало пион--пионное рассеяние. Два
пиона могут объединиться в ро-подобный резонанс, а затем вновь появиться как два
пиона. На обычной диаграмме Фейнмана это процесс аннигиляции и
повторного рождения, который принято называть процессом в \(s\)-канале. Если повернуть ту
же элементарную диаграмму набок, получится перекрёстный, или \(t\)-канальный, процесс,
в котором между пионами обменивается ро-подобная частица.

Ро-мезон не одинок. Он принадлежит реггевскому семейству, так что, по-видимому,
согласованность требует учитывать его и всех его возбуждённых партнёров в любом из
каналов. Схематически,

\begin{equation}
  \mathcal A(s,t)\simeq
  \sum_n \mathcal A^{(s)}_{\rho_n}(s,t),
  \qquad
  \mathcal A(s,t)\simeq
  \sum_n \mathcal A^{(t)}_{\rho_n}(s,t).
  \label{eq:l01-channel-sums}
\end{equation}

Эти выражения не являются двумя приближениями, которые следовало бы просто сложить.
Численные исследования показали, что суммирование башни резонансов в
канале аннигиляции уже давало приблизительно тот вклад, которого ожидали
от канала обмена. И наоборот, обменная башня могла воспроизвести
резонансное описание. Поэтому добавление обеих полных башен выглядело как
двойной счёт одной и той же физики. Обычная пертурбативная бухгалтерия требует,
чтобы, если существуют две диаграммы, были включены обе. Адронные амплитуды, напротив,
указывали на то, что описания в двух каналах были дуальны.

\begin{figure}[H]
  \centering
  \includegraphics[width=0.82\textwidth]{lecture_01_figure_01.png}
  \caption{На доске график Chew--Frautschi соединён с диаграммами аннигиляции и
  перекрёстного канала, так что спектральные подсказки и подсказки рассеяния становятся частями
  одной задачи.}
  \lectureframe{00:14:20}
  \label{fig:l01-regge-duality-board}
\end{figure}

Те же соотношения можно выделить в ясную схему. Первая панель показывает
линейный спектр; две другие представляют альтернативные разложения одной и той же
амплитуды по каналам.

\begin{figure}[H]
  \centering
  \resizebox{\textwidth}{!}{%
  \begin{tikzpicture}[>=Latex,thick]
    \begin{scope}
      \draw[->] (0,0) -- (4.2,0) node[right] {$M^2$};
      \draw[->] (0,0) -- (0,3.4) node[above] {$J$};
      \draw[blue!65!black] (0.45,0.35) -- (3.55,2.95);
      \draw[orange!80!black] (0.75,0.20) -- (3.85,2.80);
      \foreach \x/\y in {0.65/0.52,1.35/1.10,2.05/1.69,2.75/2.27,3.45/2.86}
        \fill[blue!65!black] (\x,\y) circle (0.07);
      \foreach \x/\y in {0.95/0.37,1.65/0.95,2.35/1.54,3.05/2.12,3.75/2.71}
        \fill[orange!80!black] (\x,\y) circle (0.07);
      \node at (2.1,-0.55) {parallel Regge families};
    \end{scope}
    \begin{scope}[xshift=5.7cm]
      \draw[->] (-1.5,1.2) -- (-0.3,0.25);
      \draw[->] (-1.5,-1.2) -- (-0.3,-0.25);
      \draw[decorate,decoration={snake,amplitude=1.5pt,segment length=7pt}]
        (0,0.25) -- (0,-0.25);
      \draw[->] (0.3,0.25) -- (1.5,1.2);
      \draw[->] (0.3,-0.25) -- (1.5,-1.2);
      \node at (0,-1.75) {$s$-channel tower};
    \end{scope}
    \begin{scope}[xshift=10.5cm]
      \draw[->] (-1.5,1.2) -- (1.5,1.2);
      \draw[->] (-1.5,-1.2) -- (1.5,-1.2);
      \draw[decorate,decoration={snake,amplitude=1.5pt,segment length=7pt}]
        (0,0.95) -- (0,-0.95);
      \node at (0,-1.75) {$t$-channel tower};
    \end{scope}
  \end{tikzpicture}}
  \caption{Параллельные реггевские семейства и два разложения по каналам,
  возникающие в адронной дуальности.}
  \label{fig:l01-regge-duality-redraw}
\end{figure}

\section{Одна мировая поверхность, два рассечения}

Естественной реакцией было поискать одну геометрическую историю, которую можно было бы
читать обоими способами. Представим пион парой кварк--антикварк и направим время
вверх. Вместо того чтобы рисовать только одномерные линии частиц, заполним
область, ограниченную кварковыми траекториями. Горизонтальное сечение этой
двумерной истории может показывать две входящие струны, которые соединяются, а позже
расходятся. Другое рассечение выявляет струну обмена между одной стороной и другой.
Две картинки фейнмановских каналов больше не являются независимыми
механизмами; это разные разложения одной поверхности.

\begin{figure}[H]
  \centering
  \includegraphics[width=0.78\textwidth]{lecture_01_figure_02.png}
  \caption{Топология кварковых линий на доске. Разные разрезы одной и той же
  двумерной истории дают резонансное и обменное прочтения.}
  \lectureframe{00:17:50}
  \label{fig:l01-worldsheet-topology-board}
\end{figure}

Область между кварком и антикварком даёт ещё одну идею. В
фиксированный момент времени сечение содержит две конечные точки и нечто, заполняющее весь
интервал между ними. Этот объект одномерен: это струна. Во время
взаимодействия она заметает двумерную поверхность.

\begin{definition}
Пространственно-временная история точечной частицы — это \emph{мировая линия}. Пространственно-временная
история струны — это \emph{мировая поверхность}. Струна в данный момент
представляет собой пространственное сечение своей мировой поверхности.
\end{definition}

\begin{figure}[H]
  \centering
  \includegraphics[width=0.78\textwidth]{lecture_01_figure_03.png}
  \caption{Поверхность на доске, заполненная между траекториями концов. Временное
  сечение пересекает её по одномерной струне.}
  \lectureframe{00:19:15}
  \label{fig:l01-worldsheet-filled-board}
\end{figure}

\begin{figure}[H]
  \centering
  \begin{tikzpicture}[>=Latex,scale=0.95]
    \fill[blue!10]
      (-2.6,-2.3) .. controls (-1.0,-1.0) and (-1.0,-0.4) .. (-0.2,0)
      .. controls (-1.0,0.5) and (-1.0,1.0) .. (-2.6,2.3)
      -- (2.6,2.3)
      .. controls (1.0,1.0) and (1.0,0.5) .. (0.2,0)
      .. controls (1.0,-0.4) and (1.0,-1.0) .. (2.6,-2.3) -- cycle;
    \draw[blue!65!black]
      (-2.6,-2.3) .. controls (-1.0,-1.0) and (-1.0,-0.4) .. (-0.2,0)
      .. controls (-1.0,0.5) and (-1.0,1.0) .. (-2.6,2.3);
    \draw[blue!65!black]
      (2.6,-2.3) .. controls (1.0,-1.0) and (1.0,-0.4) .. (0.2,0)
      .. controls (1.0,0.5) and (1.0,1.0) .. (2.6,2.3);
    \draw[red!75!black,dashed] (-2.9,-1.15) -- (2.9,-1.15)
      node[right] {time slice};
    \draw[orange!80!black,dashed] (-1.55,-2.5) -- (1.55,2.5)
      node[above] {crossed slicing};
    \draw[->] (3.45,-1.7) -- (3.45,1.8) node[above] {time};
    \node[blue!65!black] at (0,-2.7) {one world sheet};
  \end{tikzpicture}
  \caption{Чистая схема мировой поверхности. Различные рассечения организуют одну и ту же
  поверхность в разные описания промежуточных каналов.}
  \label{fig:l01-worldsheet-redraw}
\end{figure}

Эта геометрическая реконструкция носила поисковый характер, а не была выводом полной
теории из первых принципов. Разные люди приходили к ранней струнной
картине разными путями. Её сила состояла в том, что один и тот же протяжённый объект мог
учесть оба наблюдения: струна может вращаться и деформироваться, а её
мировая поверхность может кодировать дуальные каналы рассеяния.

\begin{classroomqa}
\audiencequestion{Почему следует вводить струноподобную перемычку, а не обычную частицу, которой обмениваются кварки?\lecturetimestamp{00:21:20}}
\lecturerresponse{В каждый данный момент имеется не просто локализованная частица обмена. Есть кварк, антикварк и структура во всех точках между ними. Эта структура может быть фундаментально непрерывной, а может быть плотным набором составляющих, приближающим континуум; первая картина этого не определяла. Важно было то, что занят весь промежуток и что наблюдаемые траектории не проявляли склонности обрываться, как обычный шнурок с грузиками на концах.}
\end{classroomqa}

\begin{classroomqa}
\audiencequestion{Может ли Стандартная модель описывать рассеяние пионов?\lecturetimestamp{00:22:45}}
\lecturerresponse{Да, хотя полезность зависит от энергетической области. При низких энергиях Стандартная модель через КХД и её низкоэнергетические описания объясняет очень многое в рассеянии пионов. При более высоких энергиях, где существенны целые башни адронных обменов, прямые вычисления становятся труднее, и струноподобная организация может оказаться полезнее.}
\end{classroomqa}

\begin{classroomqa}
\audiencequestion{Установили ли ранние эксперименты, что траектория Редже никогда не заканчивается?\lecturetimestamp{00:23:37}}
\lecturerresponse{Нет. Они установили лишь то, что наблюдаемая последовательность ещё не показывала конечной точки. Вывод о неограниченной растяжимости был наводящей на мысль экстраполяцией, а не измеренной бесконечностью.}
\end{classroomqa}

\section{Трубки потока и смена масштаба}

Современная КХД даёт старой картине теоретико-полевую интерпретацию. Кварк и антикварк выступают как противоположные источники глюонного поля. В обычной электродинамике электрический поток расплывается в стороны по мере раздвижения противоположных зарядов. Поле ослабевает с расстоянием, поэтому область между зарядами не сохраняет фиксированного поперечного профиля. КХД нелинейна. Её динамика может сконфинировать соответствующий хромоэлектрический поток в узкой трубке, соединяющей кварк и антикварк.

\begin{figure}[H]
  \centering
  \includegraphics[width=0.78\textwidth]{lecture_01_figure_04.png}
  \caption{Силовые линии поля кварка и антикварка, изображённые как узкая трубка потока, а не как свободно расплывающееся кулоновское поле.}
  \lectureframe{00:25:15}
  \label{fig:l01-qcd-flux-board}
\end{figure}

\begin{figure}[H]
  \centering
  \begin{tikzpicture}[>=Latex]
    \fill[orange!14,rounded corners=8pt] (-3.0,-0.55) rectangle (3.0,0.55);
    \draw[orange!80!black,thick,rounded corners=8pt]
      (-3.0,-0.55) rectangle (3.0,0.55);
    \fill[blue!70!black] (-3.35,0) circle (0.17);
    \fill[red!75!black] (3.35,0) circle (0.17);
    \node[below] at (-3.35,-0.22) {$q$};
    \node[below] at (3.35,-0.22) {$\bar q$};
    \foreach \y in {-0.30,0,0.30}
      \draw[->,orange!85!black] (-3.05,\y) -- (3.05,\y);
    \draw[<->] (-2.9,-1.05) -- (2.9,-1.05) node[midway,below] {$L$};
    \node at (0,0.92) {approximately fixed transverse profile};
  \end{tikzpicture}
  \caption{Сконфинированная трубка потока удлиняется, не расплываясь. Фиксированная энергия на единицу длины даёт \(E=TL\).}
  \label{fig:l01-qcd-flux-redraw}
\end{figure}

Язык частиц и язык полей по-разному выражают одно и то же. В полевом описании трубка сохраняет примерно ту же напряжённость поля и ту же площадь поперечного сечения по мере того, как её концы раздвигаются. В картине в терминах составляющих энергия растяжения создаёт между концами больше глюонной структуры. У обычной резинки число молекул фиксировано, поэтому её молекулы расходятся, пока материал не порвётся. Представьте вместо этого, что всякий раз, когда открывается промежуток, вставляется новая молекула. Плотность вдоль неё могла бы оставаться фиксированной, пока она становилась длиннее. Именно такова задуманная аналогия для глюонного поля: трубка увеличивается в длину, а не истончается по мере приближения к разрыву.

\begin{classroomqa}
\audiencequestion{Была бы струна, образующая электрон, тоже составлена из глюонов?\lecturetimestamp{00:28:08}}
\lecturerresponse{Такой вывод не следует. Адронная струна — это объект КХД на масштабе протона. Гипотетическая фундаментальная струна, имеющая отношение к квантовой гравитации, существовала бы на несравнимо меньшем планковском масштабе, и здесь нет никаких свидетельств того, что электрон является адронной глюонной струной. Поразителен тот факт, что на обоих масштабах возникла родственная математика, а не то, что физические объекты тождественны.}
\end{classroomqa}

Попытки превратить адронную картину в полную теорию взаимодействующих струн были многообещающими, но так и не дали желаемой теории адронов. Математическая модель раз за разом содержала безмассовую частицу со спином 2. В адронной теории это состояние выглядело нежелательным предсказанием; в теории гравитации у него были в точности квантовые числа гравитона. Joel Scherk и John Schwarz предложили решающую переинтерпретацию: возможно, теория говорит нам о квантовой гравитации, а не об адронах. Адронные и фундаментальные приложения физически различны, хотя именно ранняя проблема адронов указала путь к этой математике.

\begin{classroomqa}
\audiencequestion{Как линии потока идут от положительного конца к отрицательному?\lecturetimestamp{00:31:46}}
\lecturerresponse{Они выходят из одного источника и входят в другой, точно так же как силовые линии магнитного поля идут от северного полюса к южному полюсу стержневого магнита. Метки противоположных концов задают ориентацию потока; у стержневого магнита точно так же есть два противоположных полюса, а не два северных.}
\end{classroomqa}

\begin{classroomqa}
\audiencequestion{Струна непрерывна или дискретна?\lecturetimestamp{00:32:45}}
\lecturerresponse{Квантовая теория делает это различие тонким. Электромагнитное поле представляется как непрерывное поле и при этом имеет дискретные кванты. Точно так же струну можно рассматривать как непрерывную, тогда как её возбуждения квантуются. В зависимости от задаваемого вопроса оба описания могут быть верны.}
\end{classroomqa}

\begin{classroomqa}
\audiencequestion{Если изменить угол пространственно-временного среза, проходящего через мировой лист, останется ли пересечение струной?\lecturetimestamp{00:34:20}}
\lecturerresponse{Да. Изменение среза одновременности есть преобразование Лоренца. Новый срез пересекает поверхность по струне, наблюдаемой в другом состоянии движения. Даже струна, центр масс которой покоится, внутренне не обязательно неподвижна; она может вибрировать. Слова ``в покое'' относятся только к её центру масс.}
\end{classroomqa}

\section{Два разложения релятивистской энергии}

Прямое релятивистское описание протяжённого вибрирующего объекта неудобно. Полезная стратегия состоит в том, чтобы найти систему отсчёта, в которой существенное движение имеет вид обычной нерелятивистской механики. Начнём с одной частицы. В нерелятивистском случае пишут

\begin{equation}
  E_{\mathrm{nr}}=\frac{\bm p^2}{2m}+B,
  \label{eq:l01-nonrel-energy}
\end{equation}

где \(B\) — внутренняя энергия или энергия связи, не зависящая от общего импульса. Для нескольких невзаимодействующих частиц суммируются и кинетические члены, и внутренние константы. Поскольку движением управляют только разности энергий, не зависящую от состояния аддитивную константу часто можно опустить.

Точное релятивистское выражение для энергии имеет вид

\begin{equation}
  E=\sqrt{\bm p^2c^2+m^2c^4}.
  \label{eq:l01-rel-energy-c}
\end{equation}

\begin{figure}[H]
  \centering
  \includegraphics[width=0.76\textwidth]{lecture_01_figure_05.png}
  \caption{Нерелятивистская кинетическая энергия и точный релятивистский квадратный корень, где степени \(c\) были исправлены в ходе обсуждения.}
  \lectureframe{00:38:45}
  \label{fig:l01-relativistic-energy-board}
\end{figure}

\begin{classroomqa}
\audiencequestion{Не должны ли в релятивистской формуле стоять \(p^2c^2\) и \(m^2c^4\), а не те степени, которые сначала были написаны на доске?\lecturetimestamp{00:39:12}}
\lecturerresponse{Да. Требование размерностной согласованности фиксирует выражение как \(E^2=p^2c^2+m^2c^4\). После получения нерелятивистского предела мы положим \(c=1\), но сохранение этих степеней здесь делает разложение наглядным.}
\end{classroomqa}

При \(p^2c^2\ll m^2c^4\) вынесем \(m^2c^4\) за скобки и воспользуемся \(\sqrt{1+\epsilon}\simeq1+\epsilon/2\):

\begin{align}
  E
  &=mc^2\sqrt{1+\frac{p^2}{m^2c^2}}
  \notag\\
  &\simeq mc^2+\frac{p^2}{2m}.
  \label{eq:l01-low-momentum-expansion}
\end{align}

\begin{figure}[H]
  \centering
  \includegraphics[width=0.76\textwidth]{lecture_01_figure_06.png}
  \caption{Разложение по малому импульсу, дающее энергию покоя плюс привычный нерелятивистский кинетический член.}
  \lectureframe{00:41:30}
  \label{fig:l01-low-momentum-board}
\end{figure}

Это приближение требует большего, чем просто медленно движущийся центр масс. Ящик может медленно дрейфовать, тогда как частицы внутри него движутся почти со скоростью света. В этом случае внутренняя динамика всё ещё релятивистская. Обычный нерелятивистский предел требует, чтобы в выбранной системе отсчёта все существенные скорости составляющих были малы.

Есть и другой полезный предел. Выполним преобразование Лоренца всей системы вдоль оси \(z\), пока каждая составляющая не будет иметь большой положительный продольный импульс \(P_{zi}\). Любая составляющая, первоначально двигавшаяся назад, при достаточно большом бусте будет преобразована в движущуюся вперёд. Поперечные импульсы и массы покоя остаются фиксированными:

\begin{equation}
  P_{zi}\gg |p_{xi}|,|p_{yi}|,m_i,
  \qquad
  p_{\perp i}^2=p_{xi}^2+p_{yi}^2,
\end{equation}

где далее считается, что \(c=1\). Разложение при больших \(P_{zi}\) даёт

\begin{align}
  E_i
  &=\sqrt{P_{zi}^2+p_{\perp i}^2+m_i^2}
  \notag\\
  &=P_{zi}\sqrt{1+\frac{p_{\perp i}^2+m_i^2}{P_{zi}^2}}
  \notag\\
  &\simeq P_{zi}
  +\frac{p_{\perp i}^2+m_i^2}{2P_{zi}}.
  \label{eq:l01-large-pz-one-particle}
\end{align}

\begin{figure}[H]
  \centering
  \includegraphics[width=0.76\textwidth]{lecture_01_figure_07.png}
  \caption{Разложение квадратного корня при больших \(P_z\), в котором поперечный импульс и масса сохраняются как малые величины.}
  \lectureframe{00:46:45}
  \label{fig:l01-large-pz-board}
\end{figure}

Для многих составляющих,

\begin{equation}
  E\simeq
  \sum_i P_{zi}
  +\sum_i\frac{p_{\perp i}^2}{2P_{zi}}
  +\sum_i\frac{m_i^2}{2P_{zi}}.
  \label{eq:l01-large-pz-sum}
\end{equation}

Первый член — это сохраняющийся полный продольный импульс. Он огромен, но при этом является фиксированным аддитивным вкладом. Сохраняем мы его или вычитаем, это не меняет ни одной разности энергий, точно так же как добавление к гамильтониану любого другого сохраняющегося числа не изменило бы уравнений движения. Поэтому динамическая часть имеет вид

\begin{equation}
  \begin{aligned}
    H_{\perp}=E-P_z
    &\simeq
    \sum_i\left(
      \frac{p_{\perp i}^2}{2P_{zi}}
      +\frac{m_i^2}{2P_{zi}}
    \right),\\
    P_z&=\sum_iP_{zi}.
  \end{aligned}
  \label{eq:l01-transverse-hamiltonian}
\end{equation}

\begin{figure}[H]
  \centering
  \includegraphics[width=0.78\textwidth]{lecture_01_figure_08.png}
  \caption{Многочастичное разложение на сохраняющийся продольный импульс, поперечную кинетическую энергию и зависящие от массы внутренние члены.}
  \lectureframe{00:49:16}
  \label{fig:l01-energy-sum-board}
\end{figure}

\section{Система отсчёта бесконечного импульса}

Предыдущий предел исторически назывался системой отсчёта бесконечного импульса, а теперь обычно формулируется в координатах светового конуса. Его смысл можно прочитать непосредственно из гамильтониана. В квантовой механике,

\begin{equation}
  i\hbar\frac{\partial}{\partial t}\ket{\psi}=H\ket{\psi}.
  \label{eq:l01-time-evolution}
\end{equation}

Все нетривиальные члены в \(H_{\perp}\) масштабируются как \(1/P_z\). По мере роста буста энергетические расщепления уменьшаются, а изменения происходят медленнее. Это и есть замедление времени. Радиоактивная система при большом бусте кажется распадающейся медленнее; это можно компенсировать, если наблюдать её пропорционально дольше. Замедление не устраняет саму динамику. Оно лишь меняет для неё масштаб времени.

\begin{figure}[H]
  \centering
  \includegraphics[width=0.78\textwidth]{lecture_01_figure_09.png}
  \caption{На доске редуцированная энергия отождествляется с гамильтонианом, порождающим замедленную эволюцию во времени системы, подвергнутой бусту.}
  \lectureframe{00:52:57}
  \label{fig:l01-hamiltonian-board}
\end{figure}

Поперечный член имеет в точности нерелятивистский вид \(p^2/(2m)\), но роль массы играет \(P_z\). Это физически разумно: чем больше продольный импульс, тем труднее отклонить частицу вбок, так что он выступает как поперечная инерция. Член \(m_i^2/(2P_{zi})\) не зависит от двумерного состояния движения и потому играет роль внутренней энергии или энергии связи. Релятивистское движение в трёх пространственных измерениях превратилось в движение, выглядящее нерелятивистским, в двух измерениях, перпендикулярных бусту.

Именно это даёт право использовать элементарную модель масс и пружин для релятивистской струны. Эта модель не делает вид, будто струна медленна в своей системе покоя. Она использует ту систему отсчёта, в которой поперечная динамика имеет требуемую нерелятивистскую структуру.

\section{Открытая струна из шариков и пружин}

Рассмотрим открытую струну, то есть струну с двумя концами. К этим концам позднее может быть что-то присоединено, но сейчас нас интересует сама струна. Аппроксимируем её шариками с поперечными положениями
\(\bm X_i(t)=(X_i(t),Y_i(t))\), соединёнными одинаковыми пружинами. Дискретный гамильтониан имеет вид

\begin{equation}
  H_N=\sum_{i=1}^{N}
  \left[
    \frac{\mu_N}{2}\dot{\bm X}_i^{\,2}
    +\frac{k_N}{2}
      \left(\bm X_{i+1}-\bm X_i\right)^2
  \right].
  \label{eq:l01-discrete-string}
\end{equation}

Один лишь кинетический член позволил бы шарикам разлететься. Именно хуковский член для ближайших соседей делает из них струну.

\begin{figure}[H]
  \centering
  \includegraphics[width=0.75\textwidth]{lecture_01_figure_10.png}
  \caption{Первый эскиз открытой струны с двумя поперечными координатами в модели шариков и пружин.}
  \lectureframe{01:00:30}
  \label{fig:l01-bead-chain-board}
\end{figure}

\begin{figure}[H]
  \centering
  \begin{tikzpicture}[>=Latex]
    \foreach \x in {0,...,8} {
      \fill[blue!65!black] ({0.85*\x},0) circle (0.09);
    }
    \foreach \x in {0,...,7} {
      \draw[decorate,decoration={zigzag,segment length=4pt,amplitude=2pt}]
        ({0.85*\x+0.10},0) -- ({0.85*\x+0.75},0);
    }
    \draw[->] (3.4,0.15) -- (3.4,1.25) node[above] {$\bm X_i(t)$};
    \draw[<->] (0,-0.65) -- (6.8,-0.65)
      node[midway,below] {$0\leq\sigma\leq\pi$};
  \end{tikzpicture}
  \caption{Дискретный регулятор для открытой струны. Масса шарика и жёсткость пружины масштабируются с шагом дискретизации до перехода к непрерывному пределу.}
  \label{fig:l01-bead-chain-redraw}
\end{figure}

\begin{figure}[H]
  \centering
  \includegraphics[width=0.78\textwidth]{lecture_01_figure_11.png}
  \caption{Дискретные кинетический и хуковский члены, записанные вместе до перехода к непрерывному пределу.}
  \lectureframe{01:01:30}
  \label{fig:l01-discrete-energy-board}
\end{figure}

По мере увеличения числа шариков масса каждого шарика \(\mu_N\) должна стремиться к нулю, чтобы вся струна сохраняла конечную инерцию. В то же время жёсткость пружины \(k_N\) должна расти. Длинную резиновую ленту легче растянуть на фиксированную величину, чем крошечный отрезок из того же материала; поэтому соседние точки, которые сближают друг с другом, должны быть соединены более жёстко, если непрерывный материал должен остаться тем же самым.

Введём вдоль струны координату \(\sigma\). Её численный диапазон задаётся по соглашению. Выбор

\begin{equation}
  0\leq\sigma\leq\pi
\end{equation}

удобен, потому что замкнутую струну позже можно будет параметризовать от \(0\) до \(2\pi\). Сумма переходит в интеграл, а разность между ближайшими соседями — в производную:

\begin{equation}
  \sum_i\longrightarrow\int_0^\pi d\sigma,
  \qquad
  \bm X_{i+1}-\bm X_i
  \longrightarrow
  \partial_\sigma\bm X.
  \label{eq:l01-continuum-replacements}
\end{equation}

После поглощения зависящих от шага констант в нормировку \(t\), \(\sigma\) и \(\bm X\) непрерывный гамильтониан принимает вид

\begin{equation}
  H_{\mathrm{string}}=
  \frac12\int_0^\pi d\sigma
  \left[
    \left(\partial_t\bm X\right)^2
    +\left(\partial_\sigma\bm X\right)^2
  \right].
  \label{eq:l01-continuum-hamiltonian}
\end{equation}

\begin{figure}[H]
  \centering
  \includegraphics[width=0.78\textwidth]{lecture_01_figure_12.png}
  \caption{Переход на доске от суммы по шарикам к \(\int_0^\pi d\sigma\), где разность ближайших соседей заменена на \(\partial_\sigma\bm X\).}
  \lectureframe{01:05:00}
  \label{fig:l01-continuum-board}
\end{figure}

Первый член — это кинетическая энергия в двух поперечных направлениях. Второй измеряет растяжение вдоль струны. Оба члена входят со знаком плюс, потому что это гамильтониан. В лагранжиане потенциальный член входил бы со знаком минус. Точные коэффициенты зависят от соглашений и на этом этапе не фиксировались; важна будет именно структура «кинетика плюс растяжение».

Отделим движение центра масс от относительного движения. Смещение каждой точки на одну и ту же величину перемещает частицу как целое. Колебания, изгиб и относительное растяжение — это внутренние движения. Протон, атом или любой другой составной объект можно называть частицей, хотя его составляющие движутся внутри. В этом отношении струна ничем не отличается.

\section{Почему гамильтониан измеряет квадрат массы}

Для медленно движущегося объекта энергия покоя пропорциональна его массе. Однако в описании с бесконечным импульсом внутренний вклад появляется как \(M^2/(2P_z)\). Как только общий продольный импульс и связанное с ним масштабирование времени зафиксированы, внутренний гамильтониан струны оказывается пропорционален квадрату инвариантной массы всего объекта:

\begin{equation}
  H_{\mathrm{int}}\propto M^2.
  \label{eq:l01-hint-mass-squared}
\end{equation}

\editorialnote{Более явно, кинематика светового конуса даёт инвариантное соотношение вида \(M^2=2P^+P^- - \bm P_\perp^2\). Поэтому в поперечной системе центра масс и при фиксированном \(P^+\) гамильтониан светового конуса \(P^-\) пропорционален \(M^2\). В лекции эти детали нормировки опущены, и сохраняется только пропорциональность, необходимая для аргумента.}

\begin{figure}[H]
  \centering
  \includegraphics[width=0.76\textwidth]{lecture_01_figure_13.png}
  \caption{На доске внутренняя энергия струны отождествляется с квадратом массы полного состояния струны.}
  \lectureframe{01:10:30}
  \label{fig:l01-mass-squared-board}
\end{figure}

Вибрирующая струна — это совокупность связанных осцилляторов. Её полное квантование откладывается, но важный спектральный факт уже виден. Энергии гармонического осциллятора растут дискретными ступенями. Поскольку эти внутренние энергии вносят вклад в \(M^2\), возбуждения струны естественным образом дают дискретные ступени в квадрате массы. Именно эту структуру и предвещали графики Редже в начале.

\begin{classroomqa}
\audiencequestion{Является ли \(M\) массой всей струны, и приведёт ли удвоение этой массы к учетверению соответствующего вклада в \(M^2\)?
\lecturetimestamp{01:12:10}}
\lecturerresponse{Да. \(M\) — это масса покоя полного состояния струны, включая колебания, растяжение и все прочие внутренние вклады. Эффективный гамильтониан в бустированной системе пропорционален \(M^2\), так что состояние с вдвое большей массой имеет вчетверо большее значение этого квадрата массы.}
\end{classroomqa}

\begin{classroomqa}
\audiencequestion{Как это описание может включать фотон, масса покоя которого равна нулю?\lecturetimestamp{01:12:52}}
\lecturerresponse{Безмассовому состоянию соответствует \(M=0\). То, каким образом положительные вклады осцилляторов могут объединяться, давая такое состояние, — важный вопрос квантованной теории, а не то, что решается приведённой здесь классической моделью. На данном этапе \(M\) просто обозначает полную массу покоя, и квантовая механика ещё не введена.}
\end{classroomqa}

Теперь равномерно растянем покоящуюся струну до физической длины \(L\). Кинетический член исчезает, и на интервале параметра длины \(\pi\),

\begin{equation}
  \left|\partial_\sigma\bm X\right|\simeq\frac{L}{\pi}.
\end{equation}

Следовательно,

\begin{align}
  H_{\mathrm{stretch}}
  &\propto
  \int_0^\pi d\sigma
    \left(\partial_\sigma\bm X\right)^2
  \notag\\
  &\simeq
  \pi\left(\frac{L}{\pi}\right)^2
  \propto L^2.
  \label{eq:l01-hooke-l2}
\end{align}

\begin{figure}[H]
  \centering
  \includegraphics[width=0.76\textwidth]{lecture_01_figure_14.png}
  \caption{Равномерно растянутая конфигурация, используемая для вывода \(H_{\mathrm{stretch}}\propto L^2\) и, следовательно, \(M\propto L\).}
  \lectureframe{01:15:30}
  \label{fig:l01-length-relation-board}
\end{figure}

Это поведение по закону Гука в эффективной поперечной механике. Но величина слева пропорциональна \(M^2\), а не \(M\). Отсюда следует, что

\begin{equation}
  M^2\propto L^2,
  \qquad
  M\propto L.
  \label{eq:l01-mass-length}
\end{equation}

В системе покоя энергия равна \(E_{\mathrm{rest}}=M\) при единицах \(c=1\). Обозначив коэффициент пропорциональности через натяжение струны \(T\),

\begin{equation}
  E_{\mathrm{rest}}=TL.
  \label{eq:l01-string-tension}
\end{equation}

Результат согласуется с картиной флюкс-трубки. Однородное электрическое или магнитное поле несёт однородную плотность энергии. Если его поток заключён в трубку фиксированного поперечного сечения, то каждая дополнительная единица длины добавляет одно и то же количество энергии. Эквивалентно, в составной картине, где новые степени свободы заполняют каждый вновь открывающийся промежуток, энергия на единицу длины фиксирована. Натяжение струны и есть эта энергия на единицу длины.

Этот вывод косвенный. Он не утверждает, что произвольная резиновая лента, разогнанная до релятивистских скоростей, автоматически приобретёт \(E=TL\). Аргумент начинается с поперечного нерелятивистского описания, следующего из большого буста, отождествляет его внутреннюю энергию с \(M^2\) и лишь затем переводит результат обратно в энергию системы покоя.

\section{Один и тот же объект в двух системах отсчёта}

Различие между \(M\) и \(M^2\) легко упустить, поэтому стоит поместить два разложения рядом. При \(c=1\),

\begin{align}
  E&=\sqrt{p^2+M^2}
  \simeq M+\frac{p^2}{2M},
  &&p\ll M,
  \label{eq:l01-rest-expansion-again}\\
  E&=\sqrt{P_z^2+M^2}
  \simeq P_z+\frac{M^2}{2P_z},
  &&P_z\gg M.
  \label{eq:l01-boost-expansion-again}
\end{align}

\begin{figure}[H]
  \centering
  \includegraphics[width=0.74\textwidth]{lecture_01_figure_15.png}
  \caption{Повторённое после перерыва разложение по малому импульсу, в котором внутренняя энергия покоя считывается как \(M\).}
  \lectureframe{01:28:00}
  \label{fig:l01-rest-expansion-board-again}
\end{figure}

\begin{figure}[H]
  \centering
  \includegraphics[width=0.74\textwidth]{lecture_01_figure_16.png}
  \caption{Дополнительное разложение по большому импульсу, в котором остаточная энергия пропорциональна \(M^2/(2P_z)\).}
  \lectureframe{01:31:00}
  \label{fig:l01-boost-expansion-board-again}
\end{figure}

Первое разложение отождествляет энергию покоя с массой, а обычную инерцию — с \(M\). Во втором сохраняющийся член \(P_z\) удаляется, поперечная инерция задаётся \(P_z\), а внутренняя энергия представлена величиной \(M^2/(2P_z)\). Масса остаётся пропорциональной длине струны в любой из двух систем отсчёта. Меняется лишь то, какая функция массы выступает в роли остаточного гамильтониана.

\begin{classroomqa}
\audiencequestion{Буст Лоренца не должен изменять физику. Почему кажется, что бустированная струна подчиняется другому закону для энергии?
\lecturetimestamp{01:25:35}}
\lecturerresponse{Нет, она не подчиняется другой физике. В обоих описаниях имеет место \(M\propto L\) и \(M^2\propto L^2\). В системе покоя соответствующая энергия равна \(M\); после большого буста сохраняющаяся величина \(P_z\) отделяется, а оставшийся гамильтониан пропорционален \(M^2\). Буст меняет удобное описание, а не инвариантное соотношение масса--длина.}
\end{classroomqa}

\section{Сверхпроводящая модель конфайнмента}

Сверхпроводник вытесняет магнитное поле из своего объема. Если до того, как материал станет сверхпроводящим, магнитный поток захвачен в достаточно узком канале, он не может растекаться в стороны и остается заключенным в трубку. Полезный мысленный эксперимент состоит в том, чтобы просверлить узкое отверстие в образце, находящемся в нормальном состоянии, вставить в него длинный стержневой магнит, охладить образец, а затем вынуть магнит. Эта конструкция схематична, а не практический лабораторный рецепт. Ее цель — выделить линию потока, у которой напряженность поля и плотность энергии приблизительно одинаковы по всей длине.

\begin{figure}[H]
  \centering
  \includegraphics[width=0.76\textwidth]{lecture_01_figure_17.png}
  \caption{Сверхпроводящая конструкция линии магнитного потока и соответствующая картина конфайнмента монополя--антимонополя.}
  \lectureframe{01:35:30}
  \label{fig:l01-superconductor-board}
\end{figure}

Предположим, в качестве еще одного мысленного эксперимента, что магнитный монополь и антимонополь можно поместить в сверхпроводящую среду и разнести. Магнитный поток между ними сожмется в трубку. Поскольку каждая единица длины трубки несет одну и ту же энергию,

\begin{equation}
  V(R)=T_{\mathrm{flux}}R.
  \label{eq:l01-linear-confinement}
\end{equation}

Сила \(F=-dV/dR\) имеет постоянную величину. Раздвигать пару становится всё более энергозатратно, поэтому для нее имеет место конфайнмент. Именно такого поведения ищут для кварков в адроне.

\begin{figure}[H]
  \centering
  \begin{tikzpicture}[>=Latex]
    \fill[blue!8,rounded corners=5pt] (-4,-1.35) rectangle (4,1.35);
    \node[blue!55!black] at (0,1.65) {superconducting medium};
    \fill[red!75!black] (-2.8,0) circle (0.17);
    \fill[blue!70!black] (2.8,0) circle (0.17);
    \node[below] at (-2.8,-0.22) {monopole};
    \node[below] at (2.8,-0.22) {antimonopole};
    \fill[orange!18,rounded corners=7pt] (-2.65,-0.42) rectangle (2.65,0.42);
    \draw[orange!85!black,thick,rounded corners=7pt]
      (-2.65,-0.42) rectangle (2.65,0.42);
    \foreach \y in {-0.22,0,0.22}
      \draw[->,orange!85!black] (-2.55,\y) -- (2.55,\y);
    \draw[<->] (-2.8,-1.00) -- (2.8,-1.00) node[midway,below] {$R$};
  \end{tikzpicture}
  \caption{Конфайнмент потока дает однородную трубку и линейный потенциал \(V(R)=T_{\mathrm{flux}}R\).}
  \label{fig:l01-superconductor-redraw}
\end{figure}

\begin{classroomqa}
\audiencequestion{Что играет роль сверхпроводящей среды в адроне?\lecturetimestamp{01:36:57}}
\lecturerresponse{Предлагаемая аналогия меняет местами электрические и магнитные роли. Обычный сверхпроводящий конденсат электрически заряжен и конфинирует магнитный поток. В картине дуального сверхпроводника КХД конденсат хромомагнитных объектов конфинирует хромоэлектрический заряд кварков. Такие объекты не появляются как изолированные частицы, поскольку в этом описании они представлены как сконденсированные в вакууме.}
\end{classroomqa}

\editorialnote{Описание дуального сверхпроводника — влиятельный механизм конфайнмента и излагается здесь в исторической терминологии лекции. Его следует читать как физическую картину вакуума КХД, а не как полный вывод конфайнмента из первых принципов.}

\begin{classroomqa}
\audiencequestion{Является ли пропорциональность энергии длине общим правилом для любой обычной струны, переведенной в сильно релятивистский режим?\lecturetimestamp{01:38:10}}
\lecturerresponse{Нет. Сначала потребовалась бы согласованная релятивистская теория этого материала, и его поведение зависело бы от деталей. Линейная зависимость энергии характерна для конфигурации, плотность которой при удлинении остается постоянной, или для потока, которому не дают распространяться в поперечном направлении. При растяжении возникают дополнительные составляющие или дополнительный равномерно заполненный полем участок длины, а не просто увеличивается расстояние между фиксированным набором молекул.}
\end{classroomqa}

\section{Поперечные степени свободы и размер частицы}

Светоконусная конструкция записывает струну непосредственно через ее поперечные координаты. Нет отдельного, свободно задаваемого продольного осциллятора, аналогичного каждому поперечному осциллятору. Такое уменьшение числа явных пространственных измерений на единицу — одна из примечательных особенностей формализма; тем не менее полная теория должна удовлетворять лоренц-инвариантности.

\begin{classroomqa}
\audiencequestion{Что происходит с остаточным движением вдоль направления \(z\)?\lecturetimestamp{01:40:10}}
\lecturerresponse{В светоконусном описании продольное движение не вводится как независимая степень свободы. Физические осцилляции описываются в поперечных направлениях, тогда как связи определяют продольную информацию. Эта экономия по духу связана с идеей о том, что гравитирующие системы могут допускать описание с одним явным измерением меньше, хотя эта связь здесь не выводится.}
\end{classroomqa}

\begin{classroomqa}
\audiencequestion{Сколько точек на самом деле известно на графиках Редже?\lecturetimestamp{01:41:15}}
\lecturerresponse{Точный максимум зависит от семейства; грубо говоря, их около 7. Экспериментально эта последовательность не бесконечна. Существенно то, что известные точки не демонстрируют ни того рода видимого разрушения, ни конечной точки разрыва, которые привычны для обычного макроскопического материала.}
\end{classroomqa}

Адронная феноменология дала историческую мотивацию, но от математической теории струн не требуется воспроизводить адроны во всех деталях. После построения формализма спрашивают, что именно он на самом деле описывает. Его безмассовое состояние со спином 2 и другие свойства естественнее указывают на гравитацию.

Кажущийся контраст между адронами и электронами можно количественно оценить с помощью энергии жесткого вращения,

\begin{equation}
  E_{\mathrm{rot}}=\frac{J^2}{2I},
  \qquad
  I\sim MR^2.
  \label{eq:l01-rotational-energy}
\end{equation}

Угловой момент квантован. Если схематически записать \(J^2\) как \(\hbar^2 n^2\), то первая вращательная ступень имеет порядок

\begin{equation}
  \Delta E_{0\to1}\sim\frac{\hbar^2}{2I}.
  \label{eq:l01-first-rotational-step}
\end{equation}

\begin{figure}[H]
  \centering
  \includegraphics[width=0.76\textwidth]{lecture_01_figure_18.png}
  \caption{Оценка первого вращательного возбуждения через момент инерции, \(\Delta E\sim\hbar^2/(2I)\).}
  \lectureframe{01:44:00}
  \label{fig:l01-rotational-board}
\end{figure}

При фиксированной массе меньший радиус означает гораздо меньший момент инерции и гораздо большую энергию возбуждения. Адроны велики в масштабе своей внутренней динамики, поэтому вращательное возбуждение обходится примерно в ГэВ и доступно в столкновениях частиц. Если у электронов есть какая-либо пространственная подструктура, то экспериментальные ограничения делают ее масштаб гораздо меньшим, и вращательные возбуждения лежали бы далеко за пределами доступных энергий. Баскетбольный мяч представляет противоположную крайность: его момент инерции настолько велик, что один квант углового момента требует незаметно малой энергии. Атомные вращательные или орбитальные ступени находятся между этими крайностями, обычно на масштабе электронвольт.

\begin{classroomqa}
\audiencequestion{Может ли кажущийся элементарным электрон всё же иметь пространственную структуру и вращательные возбуждения?\lecturetimestamp{01:42:04}}
\lecturerresponse{Скрыть их можно было бы только при очень высоких энергиях. Цена одной единицы углового момента масштабируется как \(1/I\), а \(I\sim MR^2\). Поэтому чрезвычайно малый радиус отодвигает первое возбуждение далеко на графике квадратов масс. Отсутствие наблюдаемых возбуждений электрона означает либо что электрон точечен, либо что любой масштаб композитности намного меньше того, что могут разрешить современные эксперименты.}
\end{classroomqa}

Начальный график Редже теперь можно интерпретировать динамически. Адроны обнаруживают доступные вращательные и колебательные состояния, потому что они протяженны. Равные шаги в их спектре квадратов масс — это именно тот признак, которого ожидают, когда осцилляторный гамильтониан отождествляется с \(M^2\). Их приблизительно линейная зависимость энергии от длины — признак трубки потока с постоянным натяжением. Спектроскопия, дуальность рассеяния, геометрия мирового листа и механика светового конуса — это четыре описания одной и той же объединяющей идеи: протяженного релятивистского объекта, естественной внутренней переменной которого является квадрат массы.
